% !TeX root = ../../20260524.tex

\section{優化}

\subsection{剪枝}

\begin{frame}{例題一}
    \begin{problem}[TOJ 657 八皇后問題]
        $8 \times 8$ 的棋盤上要放八隻皇后，有些位置不能放皇后

        問有多少種放法可以放下八隻皇后，並且彼此攻擊不到彼此
    \end{problem}
\end{frame}

\begin{frame}{例題一解答}
    發現總共有 $64 \times (64-1) \times \cdots \times (64-7)$ 種放置方法

    這樣太多了，有更好的枚舉方法嗎?\pause

    我們發現有些放法其實根本不可能放得下

    對於兩隻皇后來講，只要他們所在行跟所在列相同，那一定是不合法的放置方法

    所以我們每次都枚舉每一列的皇后要放在哪一行，並且第 $i$ 列的皇后不能與 $1 \sim i-1$ 列的皇后同一行

    這樣複雜度就被壓到 $O(8! \times \text{ 檢查斜線時間})$
\end{frame}

\begin{frame}{題目}
    \begin{exercise}[CSES 1625 Grid Path Description]
        有一個 $7 \times 7$ 的格子，要從左上走到左下，且將所有格子都走過(不重複走)

        那麼路徑長度一定是 $48$

        今天給你長度為 $48$ 的字串，包含 (L, R, U, D, ?)

        你可以將 ? 視為任意方向，問有多少個走法符合字串的走法，且可以從左上走到左下，並且將所有格子都走過
    \end{exercise}
\end{frame}

\subsection{Meet in the Middle}

\begin{frame}{例題二}
    \begin{problem}[CSES 1628 Meet in the Middle]
        有 $n$ 個元素 $t_i$

        問有多少個子集合總和是 $x$

        $1 \le n \le 40$
    \end{problem}
\end{frame}

\begin{frame}{例題二解答}
    枚舉子集合的做法，$O(2^n) = O(2^{40})$，爆掉了

    有沒有更快的做法?\pause

    發現枚舉前 $1 \sim \lfloor \frac{n}{2} \rfloor$ 個子集合

    將每個子集合的總和記錄下來

    再枚舉 $\lfloor \frac{n}{2} \rfloor \sim n$ 的子集合

    這樣右邊子集合的總和 $a$ 可以用二分搜找到左邊的子集合是否存在總合為 $k-a$ 的\pause

    時間複雜度降到 $2^{\lfloor \frac{n}{2} \rfloor} \times \log (\lfloor \frac{n}{2} \rfloor)$
\end{frame}